% =====================================================================
%  Chapter 7: Spatial Equilibrium and Urban Structure
%  Source: SpatialStats.tex (intro prose) + new framework
%  NOTE: This chapter requires the most new writing. Sections marked
%        % TODO are high-priority development targets.
% =====================================================================

\chapter{Spatial Equilibrium and Urban Structure}
\label{chap_SpatialEquil}

Quantitative analysis of spatial data has advanced greatly in recent years,
largely as the result of the rapidly decreasing cost of computational power,
the increased availability of high-resolution data, and methodological
innovation in spatial econometrics. Indeed, recent analytical and
technological advances are increasingly blurring the traditional boundaries
in spatial data analysis where GIS\index{geographic information systems},
remote sensing\index{remote sensing} and spatial statistics\index{spatial
statistics} were essentially treated as separate fields
\citep[e.g.][]{Davis2002a,O'Sullivan2003,Campbell2006}.

Nowhere in the social sciences, perhaps, is this silent revolution felt more
than in urban and regional planning. While planners have usually emphasised
GIS applications for land-use and environmental planning, the discipline has
never embraced all three sub-fields of quantitative spatial data analysis in
a systematic way. With the rapid deployment of ``Big Data''\index{Big Data}
and the emergence of ``policy informatics''\index{policy informatics} as a new
paradigm among computational social scientists, planners should increasingly
see themselves at the forefront of integrated spatial data methods
\citep{Anselin2012,Barrett2011}.

\section{Why Spatial Equilibrium?}
\label{sec_WhySpatEq}

The \textbf{spatial equilibrium}\index{spatial equilibrium} framework provides
the theoretical spine that connects urban economics, regional science,
economic geography, and real estate into a single coherent argument. The
central insight, due to \citet{Rosen1979} and \citet{Roback1982}, is that
workers and firms sort across locations until utility and profits are equalised
at the margin. In equilibrium, observed geographic patterns of wages, rents,
and employment reflect underlying fundamentals---productivity differences,
amenity values, and the costs of transportation and housing.

Three implications follow:

\begin{enumerate}
  \item \textbf{Wages, rents, and quality of life are jointly determined.}
    High-amenity cities will have lower wages and higher rents, all else
    equal, because workers are willing to accept lower pay to live there.
    Observed wage or rent levels alone are not informative about the
    underlying distribution of productivity and amenity.

  \item \textbf{Observed patterns are equilibrium outcomes.} If we observe
    a spatial concentration of activity, it is because the benefits of
    concentration (agglomeration economies) outweigh the costs (congestion,
    land prices). Spatial econometrics tests whether and how adjustment
    toward equilibrium propagates across space.

  \item \textbf{Policy interventions change the equilibrium.} A place-based
    subsidy shifts the equilibrium allocation of workers and firms. The
    incidence of the subsidy---who captures the benefit---depends on the
    elasticities of housing supply and labour demand. Identifying these
    elasticities is a central task of regional policy informatics.
\end{enumerate}

% TODO: Add Rosen-Roback model derivation (wage-rent gradient)


\section{The Urban Land Market}
\label{sec_LandMarket}

\subsection{The von Thünen Model and Bid-Rent Gradients}

The earliest formal spatial equilibrium model is due to \citet{VonThunen1826}.
In von Thünen's model, a city (central market) is surrounded by an
agricultural hinterland. Farmers bid for land based on the profit from
agricultural production after transport costs. Land closer to the city
commands higher rents because transport costs are lower; the equilibrium
rent gradient is downward-sloping from the centre.

The generalisation to urban land is immediate: replace agricultural produce
with commuting costs. Workers bid for residential land based on their surplus
income after paying rent and commuting costs. The equilibrium bid-rent
function for land at distance $d$ from the city centre is:
\begin{equation}
  r(d) = r_0 - t \cdot d
\end{equation}
where $r_0$ is the rent at the city centre and $t$ is the commuting cost per
unit distance.

\subsection{The Monocentric City: Alonso-Muth-Mills}

The \textbf{Alonso-Muth-Mills model}\index{monocentric city model} formalises
the urban land market equilibrium with utility-maximising households,
profit-maximising developers, and a competitive land market
\citep{Alonso1964, Muth1969, Mills1967}. In equilibrium, the spatial
gradient of house prices compensates for the gradient in commuting costs;
household utility is equalised across all locations.

% TODO: Full AMM derivation with utility function and equilibrium conditions


\section{Agglomeration and the Concentration of Economic Activity}
\label{sec_Agglomeration}

\subsection{Why Do Cities Exist?}

The existence of cities requires an explanation: why do workers and firms
pay the premium of dense urban land when they could spread out? Three
categories of agglomeration economies\index{agglomeration economies} explain
the productivity premium of density \citep{Duranton2004}:

\begin{enumerate}
  \item \textbf{Sharing:} Dense markets allow firms to share indivisible
    inputs (specialised suppliers, infrastructure) and workers to share a
    thick labour market with lower search costs.
  \item \textbf{Matching:} Thick markets improve the quality of
    worker-firm matches, increasing productivity and wages. Cities are
    matching engines \citep{Helsley1990}.
  \item \textbf{Learning:} Knowledge spillovers diffuse more readily over
    short distances. Workers in dense environments learn from each other
    through informal contact \citep{Glaeser1992}.
\end{enumerate}

\subsection{The New Economic Geography}

\citet{Krugman1991} formalised the conditions under which economic activity
concentrates in a core region rather than distributing uniformly across space.
In Krugman's model, increasing returns to scale in manufacturing create
incentives to concentrate production; forward and backward linkages between
firms and workers amplify the initial concentration. The outcome depends on
the balance between \emph{centripetal forces} (scale economies, linkages)
and \emph{centrifugal forces} (high land prices, congestion, iceberg transport
costs).

% TODO: Add NEG bifurcation diagram and formal model sketch


\section{The Housing Market as a Spatial Equilibrium Price Surface}
\label{sec_HousingMarket}

\subsection{Hedonic Pricing: The Lancaster-Rosen Model}

A property is a bundle of characteristics---location, size, quality,
neighbourhood attributes, school quality. The \textbf{hedonic regression}
\index{hedonic regression} estimates the implicit price of each characteristic:
\begin{equation}
  \ln P_i = X_i'\beta + \varepsilon_i
\end{equation}
where $P_i$ is the sale price of property $i$ and $X_i$ is a vector of
property and neighbourhood characteristics \citep{Lancaster1966, Rosen1974}.
The hedonic coefficient $\beta_k$ measures the percentage change in price
associated with a unit change in characteristic $k$---the \emph{implicit
price} of that characteristic.

\noindent\textbf{Identification challenge.} The implicit price of an amenity
(e.g.\ proximity to a park, school quality) may be confounded by unobserved
neighbourhood quality if high-amenity neighbourhoods systematically attract
higher-income residents. IV strategies---using boundary discontinuities
or policy changes as instruments---are required for causal identification
\citep{Black1999}.

\subsection{Repeat-Sales Indices}

The \textbf{repeat-sales method}\index{repeat-sales index} \citep{Bailey1963,
Case1989} estimates price change from properties that sold more than once,
controlling for unobserved time-invariant quality by taking first differences:
\begin{equation}
  \ln P_{it} - \ln P_{is} = \sum_{\tau=s+1}^{t} \gamma_\tau D_\tau + \varepsilon_{it}
\end{equation}
where $D_\tau$ are time dummies and $\gamma_\tau$ traces the aggregate price
index. This is the methodology underlying the Case-Shiller index and the
FHFA House Price Index.

% TODO: Add worked example with Zillow ZHVI data in R


\section{Labour Market Geography}
\label{sec_LabourGeo}

\subsection{Commuting Zones and Local Labour Markets}

\citet{Autor2013} delineate \textbf{commuting zones}\index{commuting zones}
as the natural unit of local labour market analysis---functional economic
areas that contain the bulk of within-labour-market commuting flows.
The LEHD/LODES data (Chapter~\ref{chap_Data}) provide the empirical basis
for constructing commuting zones at fine geographic resolution using
workplace-to-residence flow data.

\subsection{The Spatial Mismatch Hypothesis}

\citet{Kain1968} argued that residential segregation combined with suburban
employment decentralisation created a geographic mismatch between where
low-income minority workers lived and where jobs were located.
Testing this hypothesis requires both a measure of job accessibility and
a credible instrument for residential location---since households sort into
locations partly on the basis of employment access. The LEHD/LODES data,
combined with the routing tools in the \texttt{osmdata} and \texttt{dodgr}
R packages, provide the inputs for constructing contemporary spatial
mismatch measures.

% TODO: Add commuting zone delineation example with LODES data

\renewcommand\bibname{Further reading}
\begin{subappendices}
\bibliographystyle{econometrica}
\bibliography{Literature/OverLord}
\IfFileExists{Literature/SpatialEquil.bbl}{\input{Literature/SpatialEquil.bbl}}{\typeout{Need bibtex on SpatialEquil}}
\end{subappendices}
